\documentclass[a4paper,12pt]{article}
\usepackage{graphicx}
\usepackage{amsmath}
\usepackage{amsfonts}
\usepackage{amssymb}
\usepackage{geometry}
\usepackage{booktabs}
\usepackage{float}
\usepackage{hyperref}
\usepackage{caption}
\usepackage{fancyhdr}
\usepackage{listings}
\usepackage{subcaption}
\usepackage{xcolor}
\usepackage{booktabs}
\usepackage[table]{xcolor}
\usepackage{tcolorbox} 
\usepackage{float}
\usepackage{graphicx}
\usepackage{subcaption}

\lstset{
    language=Python,            % Specify the programming language
    basicstyle=\ttfamily,       % Basic font style
    keywordstyle=\color{blue},  % Style for keywords
    commentstyle=\color{green}, % Style for comments
    stringstyle=\color{red},    % Style for strings
    frame=single,               % Add a frame around the code
    numbers=left,               % Line numbers on the left
    numberstyle=\tiny\color{gray}, % Style for line numbers
    breaklines=true,            % Allow line breaking
}

\geometry{margin=1in}
\hypersetup{
    colorlinks=true,
    linkcolor=blue,
    filecolor=magenta,      
    urlcolor=cyan,
}
\pagestyle{fancy}
\fancyhf{}
\fancyhead[L]{Detection and Tracking}
\fancyhead[R]{\thepage}

\begin{document}

\begin{titlepage}
    \centering
    {\large \textbf{Project on Signals and Systems}} \\[1.5cm]

    \begin{figure}[ht]
        \centering
        \includegraphics[width=0.4\textwidth]{C:/Users/amira/Desktop/mainproj/Signal_Project_2025/latex/images/logo.jpg}
    \end{figure}

    \textbf{Department of Electrical Engineering} \\ 
    Sharif University of Technology \\[0.5cm]
    \textbf{Instructor: Dr. Mohamadmehdi Mojahedian} \\[1.5cm]
    \textbf{Submitted by:} \\[0.5cm]
    Amirali Yazdiani - 402102765 \\[0.5cm]
    Saba Sheikh Esmaili - 402102002 \\[0.5cm] 
    Shayan Karimi - 402102335 \\[2cm]
    \textbf{Date:} \\[0.5cm]
    July 2025
\end{titlepage}


\newpage

\tableofcontents
\newpage

\setcounter{secnumdepth}{1} 

\section{Introduction and Overview}

This project is structured around three core components, each playing a vital role in the overall system:

\begin{enumerate}
    \item \textbf{Object Detection} – Identifying and locating objects within input data using advanced computer vision techniques.
    \item \textbf{Object Tracking} – Continuously monitoring the detected objects across frames to understand their motion and behavior.
    \item \textbf{Final Product} – Integrating the detection and tracking modules into a complete, functional system ready for deployment.
\end{enumerate}

In the sections that follow, we will provide a comprehensive explanation of each component, detailing the methods used, challenges encountered, and results obtained.

\section{Object Detection}

There are various techniques available for detecting objects within an image. In this project, we focus on two widely used and powerful methods:

\begin{itemize}
    \item \textbf{YOLO (You Only Look Once)}
    \item \textbf{Faster R-CNN (Region-based Convolutional Neural Network)}
\end{itemize}

In the following sections, we will provide a detailed explanation of how each of these methods works, highlighting their architectures, key differences, and practical examples.

\subsection{YOLO}

YOLO is a real-time object detection method, which is essential for our project. As the name suggests, this approach processes the image only once — hence "You Only Look Once." It is known for being both fast and efficient.

Given an input image, YOLO provides:
\begin{itemize}
    \item The objects present in the image
    \item The locations of those objects (bounding boxes)
    \item The class each object belongs to
\end{itemize}

\subsubsection{How it Works}
\begin{enumerate}
    \item \textbf{Dividing} \\ 
    The image is divided into an \(S \times S\) grid.
    
    \item \textbf{Bounding Box Prediction} \\ 
    Each grid cell predicts bounding boxes with two key parameters: the coordinates and a confidence score.
    
    \item \textbf{Class Prediction} \\ 
    Each cell also predicts the probability distribution over possible object classes.
    
    \item \textbf{Filtering} \\ 
    Using the confidence score and class probabilities, the model determines which objects are present and removes overlapping predictions using non-maximum suppression.
\end{enumerate}

\subsubsection{Implementation Examples}

You can access the code at \href{https://github.com/amirali83/Signal_Project_2025/tree/main/Codes/Detection}{this GitHub repository}.
These are some examples using code above.

\begin{figure}[H]
    \centering
    \begin{subfigure}[b]{0.45\textwidth}
        \includegraphics[width=\textwidth]{C:/Users/amira/Desktop/mainproj/Signal_Project_2025/latex/images/Detection/person1.jpg}
        \caption{raw image}
        \label{fig:image1}
    \end{subfigure}
    \hfill
    \begin{subfigure}[b]{0.45\textwidth}
        \includegraphics[width=\textwidth]{C:/Users/amira/Desktop/mainproj/Signal_Project_2025/latex/images/Detection/person1Yolo.jpg}
        \caption{processed image}
        \label{fig:image2}
    \end{subfigure}
    \caption{person 1 video}
    \label{fig:side_by_side}
\end{figure}

\begin{figure}[H]
    \centering
    \begin{subfigure}[b]{0.45\textwidth}
        \includegraphics[width=\textwidth]{C:/Users/amira/Desktop/mainproj/Signal_Project_2025/latex/images/Detection/person2.jpg}
        \caption{raw image}
        \label{fig:image1}
    \end{subfigure}
    \hfill
    \begin{subfigure}[b]{0.45\textwidth}
        \includegraphics[width=\textwidth]{C:/Users/amira/Desktop/mainproj/Signal_Project_2025/latex/images/Detection/person2Yolo.jpg}
        \caption{processed image}
        \label{fig:image2}
    \end{subfigure}
    \caption{person 2 video}
    \label{fig:side_by_side}
\end{figure}

A webcam can also be used to show that the process runs in real time.

\begin{figure}[H]
    \centering
    \begin{subfigure}[b]{0.45\textwidth}
        \includegraphics[width=\textwidth]{C:/Users/amira/Desktop/mainproj/Signal_Project_2025/latex/images/Detection/real.jpg}
        \caption{raw image}
        \label{fig:image1}
    \end{subfigure}
    \hfill
    \begin{subfigure}[b]{0.45\textwidth}
        \includegraphics[width=\textwidth]{C:/Users/amira/Desktop/mainproj/Signal_Project_2025/latex/images/Detection/realYolo.jpg}
        \caption{processed image}
        \label{fig:image2}
    \end{subfigure}
    \caption{real time process}
    \label{fig:side_by_side}
\end{figure}

\subsection{Faster R-CNN}

Faster R-CNN stands for "Faster Region-based Convolutional Neural Network." It’s one of the most accurate and widely-used object detection algorithms.
It is based on R-CNN and Fast R-CNN and also using RPN (Region Proposal Network) which makes it faster.

\subsubsection{How it Works}

\begin{enumerate}
    \item Extracting Features \\
        Image is passed through a CNN backbone to extract features.
    \item RPN \\
        From the extracted features RPN try to propose object regions.
    \item ROI Pooling \\
        Object regions are resized so they can be processed again.
    \item Classification and Boxing \\
        Regions are classified and box sizes are refined.
\end{enumerate}

\subsubsection{Implementation Examples}

You can access the code at \href{https://github.com/amirali83/Signal_Project_2025/tree/main/Codes/Detection}{this GitHub repository}.
These are some examples using code above.

\begin{figure}[H]
    \centering
    \begin{subfigure}[b]{0.45\textwidth}
        \includegraphics[width=\textwidth]{C:/Users/amira/Desktop/mainproj/Signal_Project_2025/latex/images/Detection/person1.jpg}
        \caption{raw image}
        \label{fig:image1}
    \end{subfigure}
    \hfill
    \begin{subfigure}[b]{0.45\textwidth}
        \includegraphics[width=\textwidth]{C:/Users/amira/Desktop/mainproj/Signal_Project_2025/latex/images/Detection/person1CNN.jpg}
        \caption{processed image}
        \label{fig:image2}
    \end{subfigure}
    \caption{person 1 video}
    \label{fig:side_by_side}
\end{figure}

\begin{figure}[H]
    \centering
    \begin{subfigure}[b]{0.45\textwidth}
        \includegraphics[width=\textwidth]{C:/Users/amira/Desktop/mainproj/Signal_Project_2025/latex/images/Detection/person2.jpg}
        \caption{raw image}
        \label{fig:image1}
    \end{subfigure}
    \hfill
    \begin{subfigure}[b]{0.45\textwidth}
        \includegraphics[width=\textwidth]{C:/Users/amira/Desktop/mainproj/Signal_Project_2025/latex/images/Detection/person2CNN.jpg}
        \caption{processed image}
        \label{fig:image2}
    \end{subfigure}
    \caption{person 2 video}
    \label{fig:side_by_side}
\end{figure}

It can also be used with a webcam; however, it significantly reduces the frame rate and causes noticeable lag on the computer 
because of heavy computing that it needs.

\subsection{Overall Conclusion}

Based on the FPS that we calculated with each method and the fact that we need to have a real-time function we are going to use YOLO for 
our detection in the final part.

\section{Object Tracking}
There are a lot of ways to detect track moving objects we are going to focus on three widely used methods.

\begin{itemize}
    \item \textbf{CSRT (Channel and Spatial Reliability Tracker)}
    \item \textbf{MOSSE (Minimum Output Sum of Squared Error tracker)}
    \item \textbf{KCF (Kernel Correlation Filter)}
\end{itemize}

\subsection{CSRT}

CSRT is an object tracking algorithm that use DCF and it emphasizes on accuracy more than anything else.
You can use for cases like:

\begin{itemize}
    \item \textbf{Rotation}
    \item \textbf{Occlusion (when object is blocked)}
    \item \textbf{Object gets bigger or smaller}
    \item \textbf{Object being similar to background}
\end{itemize}

\subsubsection{How it Works}

There are four main steps:
\begin{enumerate}
    \item \textbf{DCF (Discriminative Correlation Filter)} \\
        Algorithm uses a correlation filter to find the difference between background and object it compares next frame with the main frame to compare differences.
    \item \textbf{Spatial Reliability Map} \\
        It tries to ignored background noise and focus on parts that are reliable. Also tries to track pixels that help most to track motion.
    \item \textbf{Channel Reliability} \\
        There are different factors in images like color, gradients, etc it will give each one of these different weights based on how reliable each one of these are.
    \item \textbf{Model Update} \\
        Based on the previous steps it will update the parameters to give better results in the next frame.
\end{enumerate}

\subsubsection{Implementation Examples}

You can access the code at \href{https://github.com/amirali83/Signal_Project_2025/tree/main/Codes/Tracking}{this GitHub repository}. \\
Also you can see the results in \href{https://github.com/amirali83/Signal_Project_2025/tree/main/Results/CSRT}{this GitHub repository}.

\subsubsection{Results}

\begin{figure}[H]
    \centering
    \begin{subfigure}[b]{0.45\textwidth}
        \centering
        \includegraphics[width=\linewidth]{C:/Users/amira/Desktop/mainproj(1)/Signal_Project_2025/Results/CSRT/FPScar1.jpg}
        \caption{Car 1}
    \end{subfigure}
    \hfill
    \begin{subfigure}[b]{0.45\textwidth}
        \centering
        \includegraphics[width=\linewidth]{C:/Users/amira/Desktop/mainproj(1)/Signal_Project_2025/Results/CSRT/FPSperson1.jpg}
        \caption{Person 1}
    \end{subfigure}

    \vspace{0.5em} % small vertical space

    \begin{subfigure}[b]{0.45\textwidth}
        \centering
        \includegraphics[width=\linewidth]{C:/Users/amira/Desktop/mainproj(1)/Signal_Project_2025/Results/CSRT/FPSperson2.jpg}
        \caption{Person 2}
    \end{subfigure}
    \hfill
    \begin{subfigure}[b]{0.45\textwidth}
        \centering
        \includegraphics[width=\linewidth]{C:/Users/amira/Desktop/mainproj(1)/Signal_Project_2025/Results/CSRT/FPSperson3.jpg}
        \caption{Person 3}
    \end{subfigure}

    \caption{Tracking results using the CSRT algorithm across different videos.}
    \label{fig:csrt-results}
\end{figure}

We can clearly observe a significant drop in FPS, with the frame rates for the four different videos being 22, 4, 6, and 8, respectively — all of which are relatively low, indicating suboptimal real-time performance.

\subsection{MOSSE}

Like CSRT, MOSSE is a tracking algorithm, but its strength lies in its speed. It can be used in the same scenarios as CSRT, but it may not perform as well in certain challenging situations and might fail under specific conditions.

\subsubsection{How it Works}
We are not going to talk about the math behind it as it is already written for us. 
The core idea of MOSSE is to find a corelation filter that when it is applied to an image it will produce a sharp peak at the targets location and low in other areas.
Also the calculation are done in frequency domain cause it will be faster than time domain.
And at last by updating the model it will try to get better results.

\subsubsection{Implementation Examples}

You can access the code at \href{https://github.com/amirali83/Signal_Project_2025/tree/main/Codes/Tracking}{this GitHub repository}. \\
Also you can see the results in \href{https://github.com/amirali83/Signal_Project_2025/tree/main/Results/MOSSE}{this GitHub repository}.

\subsubsection{Results} 

\begin{figure}[H]
    \centering
    \begin{subfigure}[b]{0.45\textwidth}
        \centering
        \includegraphics[width=\linewidth]{C:/Users/amira/Desktop/mainproj(1)/Signal_Project_2025/Results/CSRT/FPScar1.jpg}
        \caption{Car 1}
    \end{subfigure}
    \hfill
    \begin{subfigure}[b]{0.45\textwidth}
        \centering
        \includegraphics[width=\linewidth]{C:/Users/amira/Desktop/mainproj(1)/Signal_Project_2025/Results/MOSSE/FPSperson1.jpg}
        \caption{Person 1}
    \end{subfigure}

    \vspace{0.5em} % small vertical space

    \begin{subfigure}[b]{0.45\textwidth}
        \centering
        \includegraphics[width=\linewidth]{C:/Users/amira/Desktop/mainproj(1)/Signal_Project_2025/Results/MOSSE/FPSperson2.jpg}
        \caption{Person 2}
    \end{subfigure}
    \hfill
    \begin{subfigure}[b]{0.45\textwidth}
        \centering
        \includegraphics[width=\linewidth]{C:/Users/amira/Desktop/mainproj(1)/Signal_Project_2025/Results/MOSSE/FPSperson3.jpg}
        \caption{Person 3}
    \end{subfigure}

    \caption{Tracking results using the MOSSE algorithm across different videos.}
    \label{fig:csrt-results}
\end{figure}

The results demonstrate both good FPS performance and acceptable tracking quality, making this approach effective and reliable overall.

\subsection{KCF}

KCF (Kernelized Correlation Filters) a real-time object tracking method the builds on methods like MOSSE but change them bu using:

\begin{itemize}
    \item \textbf{Non-linear kernels to improve accuracy}
    \item \textbf{Circulant matrices and Fourier transforms to make it faster}
\end{itemize}

\subsubsection{How it Works}

Like last part we are not going to focus on the math behind it but key points.

\begin{enumerate}
    \item \textbf{Correlation Filters for Tracking} \\
        Like MOSSE tries to pin a location for the target.
    \item \textbf{Circulant Matrix Trick} \\
        This method makes a circulant matrix by saving all cyclic shifts of an image without storing them and by using FFT it will compute responses efficient.
    \item \textbf{Kernel Trick}
\end{enumerate}

And finally by updating model it will try to get better.
\subsubsection{Implementation Examples}

You can access the code at \href{https://github.com/amirali83/Signal_Project_2025/tree/main/Codes/Tracking}{this GitHub repository}. \\
Also you can see the results in \href{https://github.com/amirali83/Signal_Project_2025/tree/main/Results/KCF}{this GitHub repository}.

\subsubsection{Results} 

\begin{figure}[H]
    \centering
    \begin{subfigure}[b]{0.45\textwidth}
        \centering
        \includegraphics[width=\linewidth]{C:/Users/amira/Desktop/mainproj(1)/Signal_Project_2025/Results/KCF/FPScar1.jpg}
        \caption{Car 1}
    \end{subfigure}
    \hfill
    \begin{subfigure}[b]{0.45\textwidth}
        \centering
        \includegraphics[width=\linewidth]{C:/Users/amira/Desktop/mainproj(1)/Signal_Project_2025/Results/KCF/FPSperson1.jpg}
        \caption{Person 1}
    \end{subfigure}

    \vspace{0.5em} % small vertical space

    \begin{subfigure}[b]{0.45\textwidth}
        \centering
        \includegraphics[width=\linewidth]{C:/Users/amira/Desktop/mainproj(1)/Signal_Project_2025/Results/KCF/FPSperson2.jpg}
        \caption{Person 2}
    \end{subfigure}
    \hfill
    \begin{subfigure}[b]{0.45\textwidth}
        \centering
        \includegraphics[width=\linewidth]{C:/Users/amira/Desktop/mainproj(1)/Signal_Project_2025/Results/KCF/FPSperson3.jpg}
        \caption{Person 3}
    \end{subfigure}

    \caption{Tracking results using the KCF algorithm across different videos.}
    \label{fig:csrt-results}
\end{figure}

Although the FPS drop is not as significant as observed in CSRT, a slight decrease in performance is still noticeable. Additionally, the tracker does not perform reliably, and several tracking inaccuracies can be observed throughout the sequences.

\subsection{Overall Conclusion}
Based on the observed FPS (frames per second) from each tracking method, and considering that the main priority is to achieve real-time performance, the CSRT algorithm is clearly not suitable for our needs due to its relatively slow speed.

This narrows our options down to KCF and MOSSE, both of which offer significantly higher speed and are capable of running in real-time.

To make it short:

\begin{itemize}
    \item MOSSE is the fastest, offering extremely high frame rates, but it may suffer from reduced accuracy and robustness in complex scenarios.
    \item KCF, while slightly slower than MOSSE, provides better tracking accuracy thanks to its use of kernelized correlation filters, though it still lacks built-in support for scale variation.
\end{itemize}



\begin{table}[H]
\centering
\renewcommand{\arraystretch}{1.3}
\begin{tabular}{@{} lccc @{}}
\toprule
\rowcolor{gray!20}
\textbf{Feature} & \textbf{MOSSE} & \textbf{KCF} & \textbf{CSRT} \\
\midrule
Tracking Speed (FPS)        & \textbf{Very High} & High           & Moderate        \\
Accuracy                    & Low–Medium         & Medium         & \textbf{High}   \\
Robustness to Occlusion     & Poor               & Fair           & \textbf{Good}   \\
Scale Adaptation            & No                 & No             & \textbf{Yes}    \\
Rotation Handling           & No                 & No             & Partial         \\
Initialization Time         & \textbf{Fast}      & \textbf{Fast}  & Slower          \\
Model Size / Memory         & \textbf{Very Low}  & Low            & \textbf{High}   \\
Supports Online Learning    & Yes                & Yes            & Yes             \\
Performance in Clutter      & Poor               & Fair           & \textbf{Good}   \\
Recovery from Loss          & No                 & No             & \textbf{Yes}    \\
Best Use Case               & Real-time speed    & Balanced       & Complex tracking \\
\bottomrule
\end{tabular}
\caption{Qualitative comparison of MOSSE, KCF, and CSRT tracking algorithms}
\label{tab:trackers}
\end{table}

\section{Our Previous Algorithm}

We have two algorithms cause in the first one we used KCF which we were not allowed to use and wev should have written an algorithm from scratch. \\
You can access the code at \href{https://github.com/amirali83/Signal_Project_2025/tree/main/Codes/Ours}{this GitHub repository}. \\
Now, let's discuss the steps taken to develop the final code.

\subsection{Step One}

We start by using YOLO to detect objects that appear in the video.
Typically, YOLO is run once at the beginning to detect objects, and then tracking algorithms maintain their locations 
across frames. However, our goal included detecting any new objects entering the scene at any time. 
To achieve this, we run YOLO every 5 frames to identify newly appeared objects. This approach allows us to continuously 
detect new entries but can sometimes create an illusion of continuous tracking because objects are re-detected frequently. 
The output of this step is a video with bounding boxes drawn around detected objects at regular intervals.


\subsection{Step Two}

We evaluated three different tracking algorithms but faced a major constraint: the system needed to run in real-time. 
Additionally, it was essential for the tracker to maintain object identities even during occlusions. 
Considering these requirements, we chose the KCF (Kernelized Correlation Filter) algorithm, as it offers a good balance between speed and robustness. 
Furthermore, our solution supports multi-object tracking and assigns a unique ID to each tracked object, ensuring consistent identification throughout the video.

\subsection{Step three}

One approach we implemented was using a Kalman Filter for motion prediction. 
This technique enabled smoother and more accurate object tracking by estimating the object's future position based on its past movement.

\subsection{Step Four}

After that are some small details like:

\begin{itemize}
    \item Dynamic Bounding Box Resizing
    \item Duplicate Detection Suppression
    \item FPS Display
\end{itemize}

\subsection{Step Five}

\subsubsection{Appearance-Based Re-Identification (ReID)}

One of the advanced features we experimented with was Appearance-Based Re-Identification (ReID). 
Although the results did not fully meet our expectations, the core idea was to maintain a consistent identity for each object across frames—even when it temporarily leaves and re-enters the scene. 
For instance, if an object is assigned ID number 4 when it first appears, ReID aims to ensure it retains the same ID when it comes back into view later.

\subsubsection{How it Works}

ReID operates by extracting distinctive appearance features from each detected object and matching them against previously seen objects. 
This is done using a pre-trained neural network—specifically, ResNet-18—as a feature extractor. The extracted features are then compared using cosine similarity to determine if a newly detected object closely resembles any of the previously tracked objects.

We could use deeper ReID but we did not want this function disturb our real-time ability.

\section{Our Algorithm}

You can access the code at \href{https://github.com/amirali83/Signal_Project_2025/tree/main/Codes/Ours}{this GitHub repository}. \\
Now lets walk through this algorithm.

\subsection{Detection}
\subsubsection{Model}

We were given two options for object detection: YOLO and Faster R-CNN. 
Several factors must be considered in making this choice. One key requirement is that the program should operate in real time, prioritizing speed. 
Faster R-CNN, while generally offering higher accuracy, is significantly slower due to its two-stage detection process. In contrast, YOLO provides sufficiently accurate results while maintaining high speed, 
making it a more suitable choice for our real-time application.
So we proceed with YOLO for this part.

When using YOLO, the model returns a list of all detected objects within a frame. Each object is associated with a bounding box, defined by coordinates $(x_1, y_1)$ and $(x_2, y_2)$, as well as a corresponding class label. 

With this output, we could proceed directly to the tracking stage. However, we extended the approach by incorporating the Shi-Tomasi corner detector to extract additional keypoints, enhancing the robustness of object tracking.

\subsubsection{Shi-Tomasi Corner Detector}

As the name suggests, this algorithm focuses on detecting corners within a frame, which can then be used for tracking. \\[0.5cm]
\textbf{Why Corners?}

There are several reasons why corners are advantageous for tracking:

\begin{itemize}
    \item Corners are relatively stable and easier to detect consistently across frames.
    \item They can be extracted more reliably from the background compared to edges or flat regions.
    \item Corners typically exhibit less motion distortion than edges or uniform areas.
    \item They can be efficiently tracked using optical flow methods (discussed further in the tracking section).
\end{itemize}

\subsubsection{How it Works}

To determine whether a pixel is a corner, we first select an $n \times n$ window (or grid) centered around the pixel.

Next, we compute the structure tensor \( M \), defined as:

\[
M = \sum_{\text{window}} 
\begin{bmatrix}
I_x^2 & I_x I_y \\
I_x I_y & I_y^2
\end{bmatrix}
\]

Here, \( I_x \) and \( I_y \) represent the image gradients in the horizontal and vertical directions, respectively. The summation is performed over all pixels in the selected window.

By computing the eigenvalues (\( \lambda_1 \) and \( \lambda_2 \)) of the matrix \( M \), 
we can determine whether a given pixel corresponds to a corner. These eigenvalues represent the intensity variation in different directions within the local window. Based on their values, there are three possible interpretations:

\begin{enumerate}
    \item Both \( \lambda_1 \) and \( \lambda_2 \) are large \( \rightarrow \) the region is a \textbf{corner}.
    \item One eigenvalue is large while the other is small \( \rightarrow \) the region is likely an \textbf{edge}.
    \item Both \( \lambda_1 \) and \( \lambda_2 \) are small \( \rightarrow \) the region is a \textbf{flat surface}.
\end{enumerate}

\[Corner score = min( \lambda_1, \lambda_2)\]

So if our corner score is bigger than a chosen number that region is a corner.

\subsection{Tracking}

\subsubsection{Optical Flow}
There are three main types of optical flow:

\begin{itemize}
    \item \textbf{Lucas-Kanade}
    \item \textbf{Horn-Schunck}
    \item \textbf{Deep Learning-based methods}
\end{itemize}

In our program, we will be using the \textbf{Lucas-Kanade} method. \\
\textbf{How It Works}

First, let's discuss what tracking means. Consider a video file \( I \), where \( I(x, y, t) \) represents the pixel at coordinates \( (x, y) \) at time \( t \).  
We say a pixel is being tracked if we can find displacements \( dx \) and \( dy \) such that:

\[
I(x, y, t) = I(x + dx, y + dy, t + 1)
\]

To solve this, we use approximations and a first-order Taylor series expansion, which leads to:

\[
I_x \cdot dx + I_y \cdot dy = -I_t
\]

Here, \( I_x \), \( I_y \), and \( I_t \) are the partial derivatives of the image with respect to \( x \), \( y \), and \( t \), respectively.

This gives us one equation with two unknowns — \( dx \) and \( dy \). To address this underdetermined system, we consider a small window (or grid) of pixels around the point of interest. For example, we use a window size of \( 15 \times 15 \).  
By applying the same equation to all pixels within this window, we obtain a solvable system of equations that allows us to estimate the motion vector \((dx, dy)\).\\ \\
\textbf{Limitations and Improvements}

There are some limitations to this approach. For instance, if a pixel undergoes a large movement between frames, the method may fail to track it accurately. Moreover, while this method is relatively fast, it can still be computationally intensive for high-resolution videos.

To address these issues, we employ an image pyramid — a technique that enables tracking at multiple scales, improving both performance and robustness in the presence of large motions.

\subsubsection{Pyramids}

In simple terms, a pyramid is a collection of images derived from a single original image, each at progressively lower resolutions.

To understand this better, consider a 2-level pyramid applied to an image \( I \) of size \( 1600 \times 800 \) pixels. We generate a new image from the original by performing two key operations:

\begin{itemize}
    \item Reduce the number of pixels in each row and column by half.
    \item Attempt to preserve as much detail as possible from the original image.
\end{itemize}

This results in an image with only \( 25\% \) of the original resolution. \\ \\
There are several ways to achieve this downsampling:

\begin{itemize}
    \item Remove all pixels located at odd or even indices.
    \item Group neighboring pixels (e.g., 2-by-2 blocks) and average their values to form a new pixel.
    \item Use specialized filtering matrices (e.g., Gaussian filters) to reduce resolution while retaining important features.
\end{itemize}

After this operation, the resulting image has a size of \( 800 \times 400 \), which we refer to as the \textbf{level 1} image. The original image is referred to as \textbf{level 0}. We can continue this process for more levels, depending on the number of pyramid levels we require. \\ \\
\textbf{How it Works}

When performing tracking using image pyramids, we start from the \textbf{highest level} of the pyramid — that is, the image with the \textbf{lowest resolution}.  
Because this image contains fewer pixels, the tracking process is computationally efficient.

Once we obtain a rough estimate of the pixel's motion at this coarse level, we move \textbf{one level down} in the pyramid — to a higher-resolution image — and \textbf{refine} our tracking result.  
Since we already have an approximate idea of where the pixel moved, this refinement step is significantly faster and more accurate than starting from scratch at that level.

We repeat this process, gradually moving down the pyramid, until we reach \textbf{level 0} (the original image).  
This coarse-to-fine approach allows us to efficiently and robustly track large motions that would otherwise be difficult to detect at full resolution.

\subsubsection{Where are we now}

Up to this point, we are able to detect and track an object successfully.  
However, there are situations where an object enters a special circumstance — for example, it moves behind another object, causing a temporary \textbf{occlusion}.

In such cases, we should not immediately discard the object from tracking. Instead, we aim to continue tracking it for as long as possible, even when it is not directly visible.

To achieve this, we introduce the \textbf{Kalman Filter} — a powerful tool for prediction in the presence of uncertainty.  

\subsubsection{Kalman Filter}

The \textbf{Kalman Filter} helps us estimate the position of an object based on its previous trajectory, even if the object is temporarily occluded or not detected in the current frame.

The Kalman Filter operates in two main stages:

\begin{itemize}
    \item \textbf{Prediction}
    \item \textbf{Correction}
\end{itemize}

\textbf{Prediction} 

In the prediction stage, the Kalman Filter estimates the object's position in the next frame based on its current state and motion. The basic prediction equation is:

\[
\hat{x}_{k|k-1} = A \cdot \hat{x}_{k-1|k-1}
\]

where:
- \( \hat{x}_{k|k-1} \) is the predicted state vector at time \( k \) given observations up to time \( k-1 \),
- \( \hat{x}_{k-1|k-1} \) is the estimated state from the previous frame,
- \( A \) is the state transition matrix.

We typically define the state vector \( x \) as:

\[
x = \begin{bmatrix}
    x \\
    y \\
    dx \\
    dy
\end{bmatrix}
\]

This includes the position \( (x, y) \) and velocity \( (dx, dy) \) of the object. \\[0.5cm]
\textbf{Correction}

After predicting the object's next position, we compare this prediction with the actual measurement (if available).  
If the object does not appear at the predicted location, the Kalman Filter updates its internal parameters to reduce future prediction errors.

This is done using the correction (or update) equation:

\[
\hat{x}_{k|k} = \hat{x}_{k|k-1} + K \cdot (z_k - H \cdot \hat{x}_{k|k-1})
\]

where:
- \( \hat{x}_{k|k} \) is the updated (corrected) state estimate at time \( k \),
- \( \hat{x}_{k|k-1} \) is the predicted state estimate,
- \( z_k \) is the actual measurement (observation) at time \( k \),
- \( H \) is the observation model matrix,
- \( K \) is the Kalman gain, which determines how much the prediction should be corrected based on the measurement.

This equation adjusts the prediction based on the difference between the predicted observation \( H \cdot \hat{x}_{k|k-1} \) and the actual measurement \( z_k \).  
Over time, this feedback mechanism allows the filter to learn and adapt, improving its accuracy — even when some frames are missing or noisy.

So, the Kalman Filter uses past information to predict where the object should appear in the next frame, which is especially useful when the object is temporarily hidden or its detection fails.

\subsubsection{Results} 

\begin{figure}[H]
    \centering
    \begin{subfigure}[b]{0.45\textwidth}
        \centering
        \includegraphics[width=\linewidth]{C:/Users/amira/Desktop/mainproj(1)/Signal_Project_2025/Results/OURS/FPScar1.jpg}
        \caption{Car 1}
    \end{subfigure}
    \hfill
    \begin{subfigure}[b]{0.45\textwidth}
        \centering
        \includegraphics[width=\linewidth]{C:/Users/amira/Desktop/mainproj(1)/Signal_Project_2025/Results/OURS/FPSperson1.jpg}
        \caption{Person 1}
    \end{subfigure}

    \vspace{0.5em} % small vertical space

    \begin{subfigure}[b]{0.45\textwidth}
        \centering
        \includegraphics[width=\linewidth]{C:/Users/amira/Desktop/mainproj(1)/Signal_Project_2025/Results/OURS/FPSperson2.jpg}
        \caption{Person 2}
    \end{subfigure}
    \hfill
    \begin{subfigure}[b]{0.45\textwidth}
        \centering
        \includegraphics[width=\linewidth]{C:/Users/amira/Desktop/mainproj(1)/Signal_Project_2025/Results/OURS/FPSperson3.jpg}
        \caption{Person 3}
    \end{subfigure}

    \caption{Tracking results using OUR algorithm across different videos.}
    \label{fig:csrt-results}
\end{figure}

The results are very promising, with the processing speed being so high that the video needs to be slowed down for normal viewing. Moreover, 
one of the key targets — achieving at least \( 80\% \) of the FPS performance of CSRT — has been successfully met.

\subsection{Code}

\begin{tcolorbox}[colback=lightgray!20, colframe=gray!80, title=Library Imports]
\begin{verbatim}
import cv2
import numpy as np
import time
from ultralytics import YOLO
\end{verbatim}
\end{tcolorbox}

These lines import essential libraries: OpenCV for computer vision, NumPy for numerical operations, time for FPS calculation, and the YOLO model for object detection.

\vspace{0.5em}
\begin{tcolorbox}[colback=lightgray!20, colframe=gray!80, title=Model and Video Initialization]
\begin{verbatim}
model = YOLO("yolov5s.pt")
source = "files/person1.mp4"
cap = cv2.VideoCapture(source)
\end{verbatim}
\end{tcolorbox}

The YOLOv5 model is loaded, the video source is set, and the video capture object is initialized to read frames.

\vspace{0.5em}
\begin{tcolorbox}[colback=lightgray!20, colframe=gray!80, title=Feature Detection Parameters]
\begin{verbatim}
feature_params = dict(maxCorners=100, qualityLevel=0.3, 
minDistance=7, blockSize=7)
lk_params = dict(winSize=(15, 15), maxLevel=2, 
criteria=(cv2.TERM_CRITERIA_EPS | cv2.TERM_CRITERIA_COUNT, 10, 0.03))
\end{verbatim}
\end{tcolorbox}

Defines parameters for Shi-Tomasi corner detection (to find good features to track) and Lucas-Kanade optical flow (to track those features across frames).

Some of these parameters include a maximum of 100 corners to be detected, with each corner required to have a minimum quality level of 0.3. 
Additionally, the minimum distance between any two detected corners is 7 pixels, calculated as \( \sqrt{(x_1 - x_2)^2 + (y_1 - y_2)^2} \), 
and the search is performed within a window of size \( 7 \times 7 \).


\vspace{0.5em}
\begin{tcolorbox}[colback=lightgray!20, colframe=gray!80, title=Kalman Filter Initialization]
\begin{verbatim}
kalman = cv2.KalmanFilter(4, 2)
kalman.measurementMatrix = np.array([[1, 0, 0, 0], [0, 1, 0, 0]], np.float32)
kalman.transitionMatrix = np.array([[1, 0, 1, 0], 
                                    [0, 1, 0, 1], 
                                    [0, 0, 1, 0], 
                                    [0, 0, 0, 1]], np.float32)
kalman.processNoiseCov = np.eye(4, dtype=np.float32) * 0.03
\end{verbatim}
\end{tcolorbox}

Sets up a Kalman filter to estimate and predict the object’s position and velocity, modeling motion and measurement noise.

\vspace{0.5em}
\begin{tcolorbox}[colback=lightgray!20, colframe=gray!80, title=Tracking Variables Initialization]
\begin{verbatim}
prev_gray = None
prev_pts = None
tracking = False
object_center = None
object_label = None
\end{verbatim}
\end{tcolorbox}

Initializes variables to keep track of previous frame data, feature points, tracking state, and detected object info.

\vspace{0.5em}
\begin{tcolorbox}[colback=lightgray!20, colframe=gray!80, title=FPS Calculation Setup]
\begin{verbatim}
prev_time = time.time()
fps = 0
\end{verbatim}
\end{tcolorbox}

Used to calculate the frames per second dynamically during video processing.

\vspace{0.5em}
\begin{tcolorbox}[colback=lightgray!20, colframe=gray!80, title=Main Loop - Frame Processing]
\begin{verbatim}
while True:
    ret, frame = cap.read()
    if not ret:
        break
\end{verbatim}
\end{tcolorbox}

Reads video frames one at a time. Exits the loop when no more frames are available.

\vspace{0.5em}
\begin{tcolorbox}[colback=lightgray!20, colframe=gray!80, title=FPS Computation]
\begin{verbatim}
current_time = time.time()
fps = 1.0 / (current_time - prev_time)
prev_time = current_time
\end{verbatim}
\end{tcolorbox}

Calculates the FPS based on the elapsed time since the last frame was processed.

\vspace{0.5em}
\begin{tcolorbox}[colback=lightgray!20, colframe=gray!80, title=Convert Frame to Grayscale]
\begin{verbatim}
gray = cv2.cvtColor(frame, cv2.COLOR_BGR2GRAY)
\end{verbatim}
\end{tcolorbox}

Converts the current frame to grayscale to simplify the feature tracking process.

\vspace{0.5em}
\begin{tcolorbox}[colback=lightgray!20, colframe=gray!80, title=Optical Flow Tracking]
\begin{verbatim}
if tracking and prev_pts is not None:
    next_pts, status, _ = cv2.calcOpticalFlowPyrLK(prev_gray, gray, 
    prev_pts, None, **lk_params)
    # Filter good points and update Kalman filter
    # Reset tracking if points lost
\end{verbatim}
\end{tcolorbox}

If tracking is active, uses Lucas-Kanade optical flow to track points across frames and correct the Kalman filter. If tracking fails, resets tracking status.

\vspace{0.5em}
\begin{tcolorbox}[colback=lightgray!20, colframe=gray!80, title=Object Detection and Initialization]
\begin{verbatim}
if not tracking:
    results = model(frame)[0]
    for det in results.boxes:
        # Extract bounding box and class
        # Detect good features inside the box to track
        # Initialize tracking and Kalman filter
        break
\end{verbatim}
\end{tcolorbox}

When not tracking, runs YOLO to detect objects. Initializes tracking on the detected object's features.

\vspace{0.5em}
\begin{tcolorbox}[colback=lightgray!20, colframe=gray!80, title=Kalman Prediction and Bounding Box Drawing]
\begin{verbatim}
predicted = kalman.predict()
pred_x, pred_y = int(predicted[0]), int(predicted[1])

# Calculate adaptive bounding box and predict center
# Draw rectangle and label
\end{verbatim}
\end{tcolorbox}

Predicts the object's next position with Kalman filter and draws an adaptive bounding box around the tracked features.

\vspace{0.5em}
\begin{tcolorbox}[colback=lightgray!20, colframe=gray!80, title=Display FPS and Frame]
\begin{verbatim}
cv2.putText(frame, f"FPS: {fps:.2f}", (10, 30), 
            cv2.FONT_HERSHEY_SIMPLEX, 0.9, (0, 255, 0), 2)
cv2.imshow("YOLO + Optical Flow + Kalman + Adaptive Box", frame)

if cv2.waitKey(1) & 0xFF == 27:
    break
\end{verbatim}
\end{tcolorbox}

Overlays the calculated FPS on the frame and displays it. Allows exiting the loop by pressing ESC.

\vspace{0.5em}
\begin{tcolorbox}[colback=lightgray!20, colframe=gray!80, title=Cleanup]
\begin{verbatim}
cap.release()
cv2.destroyAllWindows()
\end{verbatim}
\end{tcolorbox}

Releases video resources and closes all OpenCV windows after processing is complete.

\section{Conclusion}

In this project, we implemented and evaluated a hybrid object tracking system that combines YOLOv5 for object detection, optical flow for motion estimation, and a Kalman filter for prediction and smoothing. The system was tested on multiple video sequences and compared against standard tracking methods such as CSRT. Results show that the proposed approach offers a good balance between accuracy and performance, with competitive FPS and robust tracking in challenging conditions. Overall, the project objectives were successfully achieved.


\end{document}
